\subsection*{Examples of Lie groups}

\begin{enumerate}
\def\labelenumi{(\arabic{enumi})}
\tightlist
\item
  Let \(\mathbf{E}\) be a complete normable space over \(\mathbf{K}\) . The mapping \((x, y) \mapsto x - y\) of \(\mathbf{E} \times \mathbf{E}\) into \(\mathbf{E}\) is continuous and linear and hence analytic.
\end{enumerate}

Hence E, with its additive group and analytic manifold structures, is a Lie group.

In particular, K is a Lie group.

\begin{enumerate}
\def\labelenumi{(\arabic{enumi})}
\setcounter{enumi}{1}
\tightlist
\item
  Let A be a complete normable unital associative algebra over K. The multiplication \((x,y)\mapsto xy\) of A × A into A is bilinear and continuous and hence analytic. Proposition 3 shows that the group A* of invertible elements of A is open in A (which also follows from General Topology, Chapter IX, §3, Proposition 13) and that A* is a Lie group.
\end{enumerate}

For example, let E be a complete normable space over K and let A = \(\mathcal{L}(\mathrm{E})\) (General Topology, Chapter IX, § 3, Proposition 5). Then A* is the automorphism group GL(E) of E. This group therefore has canonically a Lie group structure over K. More particularly, GL(n, K), with the manifold structure induced by that on \(\mathbf{M}_{n}(\mathbf{K})\) , is a Lie group. For \(n = 1\) , we see that the multiplicative group K* is a Lie group with the manifold structure induced by that on K.

\begin{enumerate}
\def\labelenumi{(\arabic{enumi})}
\setcounter{enumi}{2}
\tightlist
\item
  Let G be a Lie group over K. Let \(K' = R\) or C or a non-discrete complete ultrametric field and \(\sigma\) an isomorphism of the valued field \(K'\) onto a valued subfield of K. Then the group G, with the K'-manifold structure obtained by restriction of scalars, is a Lie group over \(K'\) , which is said to be derived from the Lie group G by restriction of scalars (from K to \(K'\) by means of \(\sigma\) ). For example, every complex Lie group has canonically a real Lie group structure. Again, with every complex Lie group G is associated a complex Lie group called the conjugate of G, derived from G by means of the automorphism \(z \mapsto \bar{z}\) of C.
\end{enumerate}

